\documentclass[11pt,a4paper]{article}
\usepackage[T1]{fontenc}
\usepackage[utf8]{inputenc}
\usepackage[english]{babel}
\usepackage{amsmath,amssymb}
\usepackage[margin=2cm]{geometry}
\setlength{\parindent}{0pt}
\setlength{\emergencystretch}{3em}
\allowdisplaybreaks

\title{About \texttt{witnesses.tex}}
\author{}
\date{}
\begin{document}
\maketitle
This note explains what the test data in \texttt{witnesses.tex} is used for. That file is input for the engine in \texttt{main.c}; it is read by the test \texttt{test\_algebraic\_structures.c} in this folder.

\begin{enumerate}
\item \textbf{Purpose.} The file holds 24 true statements whose proof has to use a hypothesis of the form ``there is $k$'', or a congruence or a divisibility whose modulus is a letter. Each must be proved, with a written proof.

\item \textbf{What it exercises.} Rule 18u of \texttt{main.c}: the witness of each such hypothesis gets a name, memberships and divisibilities are replaced by their definition, and the number the thesis asks for is found as a polynomial in the named witnesses, such that the equation of the thesis is a combination of the hypotheses (rule 18s, with the extension to a factor that is never zero). The last statements have no hypothesis: the value is found when solving the equation would divide by a number that can be zero, or when two values are asked for at once.

\item \textbf{Variety.} Congruence modulo a letter is reflexive, symmetric, transitive, compatible with sum, product and square; divisibility chains; ``the difference is an integer'' among the rationals, written with $\exists$ and with $\in \mathbb{Z}$; cancellation with a B\'ezout identity; an equivalence; a negated divisibility among the hypotheses; two hypotheses and a thesis that all use the same letter $k$, as books write them.

\item \textbf{Expected result.} The line \texttt{PASS: 24 of 24 statements proved, each with a written proof}.

\item \textbf{How to run it.} From the top folder of the repository:
\begin{verbatim}
cc -std=c11 -O2 -Wall -Wextra -o /tmp/test_algebraic_structures \
    tests/algebraic_structures/test_algebraic_structures.c \
    && /tmp/test_algebraic_structures
\end{verbatim}
\end{enumerate}
\end{document}
